% Part: incompleteness
% Chapter: theories-computability
% Section: introduction

\documentclass[../../../include/open-logic-section]{subfiles}

\begin{document}

\olfileid{inc}{tcp}{int}

\olsection{ভূমিকা}

\begin{editorial}
এই অংশটি নতুন করে লেখা দরকার।
\end{editorial}

আমাদের হাতে নিচের ফলগুলি আছে:
\begin{enumerate}
\item কোনো অপেক্ষককে $\Th{Q}$-তে উপস্থাপনযোগ্য বলার অর্থের একটি সংজ্ঞা
  (\olref[req][int]{defn:representable-fn})
\item কোনো সম্বন্ধকে $\Th{Q}$-তে উপস্থাপনযোগ্য বলার অর্থের একটি সংজ্ঞা
  (\olref[req][rel]{defn:representing-relations})
\item $\Th{Q}$-এর উপস্থাপনযোগ্য অপেক্ষকগুলি ঠিক গণনসাধ্য অপেক্ষকগুলিই---এই
  বক্তব্যের একটি উপপাদ্য
  (\olref[req][int]{thm:representable-iff-comp})
\item $\Th{Q}$-এর উপস্থাপনযোগ্য সম্বন্ধগুলি ঠিক গণনসাধ্য সম্বন্ধগুলিই---এই
  বক্তব্যের একটি উপপাদ্য
  (\olref[req][rel]{thm:representing-rels})
\end{enumerate}

কোনো \emph{তত্ত্ব} হলো অনুগমনের অধীনে আবদ্ধ বাক্যসমষ্টি; অর্থাৎ, যখনই
$T$ থেকে $!A$ প্রমাণ করা যায়, তখন $!A$-ও $T$-তে থাকে। তত্ত্বকে সম্ভবত
এমন এক বাক্যসংগ্রহ হিসেবে ভাবাই ভালো, যার সঙ্গে ওই বাক্যগুলি থেকে অনুগত
সবকিছুও অন্তর্ভুক্ত। এখন থেকে $\Th{Q}$ দিয়ে আমরা
\olref[req][int]{sec}-এর আটটি স্বতঃসিদ্ধ থেকে নিষ্পাদনযোগ্য বাক্যগুলির
তত্ত্বটিকে বোঝাব। মনে রাখো, $\Th{Q}$-এর সূত্রগুলিকে সংখ্যা দিয়ে
সংকেতায়িত করা যায়; $!A$ এমন একটি সূত্র হলে $!A$-কে সংকেতায়িত সংখ্যাটি
$\Gn{!A}$ দিয়ে নির্দেশ করি। এই সংকেতায়ন ধরে এখন আমরা জিজ্ঞাসা করতে পারি,
বিভিন্ন সূত্রসমষ্টি গণনসাধ্য কি না।

\end{document}
